% !TEX root = ../main.tex

\chapter{结论与展望}
\label{chap:conclusion}

本报告整合了课程的仿真工作，按功率流向完成了从 10~kV 交流母线到等离子体负载的建模与校核，并把各环节接入了同一电气网络。主要结论如下。

\begin{enumerate}
  \item \textbf{10~kV 母线。}六个工况全部通过相量校核，默认工况母线 0.987803~pu、内部电势 1.017047~pu；重载工况在 0.110525~s 触及 1.15~pu 上限，积分状态被冻结，母线停在 0.878646~pu。
  \item \textbf{CS 整流供电。}460~s 全过程电流峰值 39.746~kA、输入功率峰值 59.619~MW、磁场 11.482~T、储能 0.7571~GJ，八个关键时刻的电流与 ODE 基准最大相差 0.247\%；六脉波详细模型验证了 300~Hz 纹波和触发顺序。
  \item \textbf{H 桥电流环。}五段给定稳态误差 0.68--1.77~A，载波 2000.0~Hz；修正死区实现后，两桥臂各 1592 次换相的间隔全部为 2~$\mu$s，无桥臂直通；调制律偏差与死区损失 17.60~V 相符，负载回路方程稳态残差均方根 2.85~V。
  \item \textbf{可变等离子体负载。}平台段约 10~kA、1~MW，磁链两条独立路径的最大残差约 0.011~Wb；线圈侧多出的功率由励磁储能增长和原边电阻损耗解释；1~$\mu$s 下电流纹波约 2.3\%，而 10~$\mu$s 步长会明显放大纹波，并使退磁尾段失去窄脉冲。
  \item \textbf{全链路联合。}整机模型连续运行 30~s，平台段母线 2199.98~V、等离子体电流 9999.28~A，门极窗口实测 2~$\mu$s 死区，退磁向交流侧净回馈 0.559~MJ。晶闸管接口的对照模型全程直流能量相对残差 0.00318\%，退磁能量由电容和制动支路吸收。两套独立模型的线圈侧能量相差 0.015\%。
  \item \textbf{接口分析。}100~MVA 母线无法支撑实验一的 59.6~MW 峰值负载，所需内部电势 1.455~pu 超出 1.15~pu 限幅，短路容量需提高到约 270~MVA 以上。
\end{enumerate}

后续工作可以沿三个方向展开：一是把联合模型的负载提升到实验一的功率等级，并按第 \ref{chap:analysis} 章的结果配置短路容量和无功补偿，观察脉冲负载对母线的暂态影响；二是对整流器 PI、电流环参数和退磁系数 $k_\psi$ 做敏感性分析；三是把等离子体负载由单端口模型扩展为带空间耦合的降阶模型，使磁链预算更接近真实放电。
